\documentclass[book.tex]{subfiles}

\begin{document}

\chapter{Imaginary Time Path Integrals}

\label{chap:imaginary_path}
\noindent\rule{11cm}{0.4pt} \\
\textbf{Prerequisites:}  \autoref{chap:imaginary_qm} \\
\textbf{Difficulty Level:} *** \\
\noindent\rule{11cm}{0.4pt}
\vspace{5mm}
\textcolor{red}{TODO: Cite Shankar}
Consider the transition probability
\begin{align}
    U(x, x'; \tau) &= \langle x | U(\tau) | x'\rangle
\end{align}
We can write this transition density
\begin{align}
    \langle x | U(\tau) | x' \rangle &= \int [\mathcal{D}x] \exp \left [ -\frac{1}{\hbar} \int_0^\tau  \mathcal{L}_E(x,\bar{x}) d\tau \right]
\end{align}
Here the path integral measure and lagrangian is given by
\begin{align}
    \int [ \mathcal{D}x] &= \lim_{N \to \infty} \left ( \frac{m}{2\pi\hbar \epsilon} \right )^{1/2}\prod_0^{N-1} \left ( \frac{m}{2\pi\hbar \epsilon}\right ) dx_i \\
    \mathcal{L}_E &= \frac{m}{2}\left ( \frac{dx}{d\tau}\right )^2 + V(x)
\end{align}
This Lagrangian is called the Euclidean lagrangian.
(\textcolor{red}{TODO: Explain how the Euclidean Lagrangian and the usual Lagrangian differ}). Similarly here the Euclidean action is given 
\end{document}